\label{ch:introduction}

\section{Purpose and how to read this document}

This document is the third inference-side companion to the \emph{Theoretical
Foundations of \texttt{tpeanuts}} document: where \emph{Inference at Daya
Bay} documents a reactor antineutrino application and \emph{Real Salt-Phase
(SNO-II)} documents a solar-neutrino application, this one documents an
\emph{atmospheric}-neutrino application, built on the only detector in this
project whose real public data is event-by-event weighted Monte Carlo
rather than a continuous cross-section fold -- fitting the atmospheric
mixing angle $\theta_{23}$ and the mass-squared splitting $\Delta m^2_{3l}$
to the real public data release of the IceCube DeepCore experiment. Every
number quoted here that is not explicitly attributed to a fitted
\texttt{tpeanuts} result is a real, published, externally verifiable
quantity, traceable through the References chapter.

The five chapters are organised as follows: this chapter introduces the
experiment itself -- geometry, detection principle, and data-taking period
-- with no reference yet to \texttt{tpeanuts}; \cref{ch:data-release}
catalogues the real official public data release the code consumes;
\cref{ch:detector-implementation} is the physics-and-mathematics core,
developing the weighted-Monte-Carlo reweighting formula and the
detector-systematics hypersurface correction \texttt{tpeanuts.detector.icecube}
implements; \cref{ch:inference-process} documents the Poisson fit machinery
and this document's own dedicated test suite and audit work;
\cref{ch:results} presents, interprets, and summarises the numerical
results obtained from \texttt{inference5\_icecube.ipynb}, compared against
IceCube's own published best fit, together with the muon-background
normalization audit this document's own working session carried out.

\section{Overview}

IceCube is a cubic-kilometre neutrino telescope instrumenting the glacial
ice at the geographic South Pole with 5160 digital optical modules (DOMs)
on 86 vertical strings, detecting the Cherenkov light emitted by charged
particles produced in neutrino interactions within or near the ice.
\textbf{DeepCore} is a denser sub-array of 8 additional strings deployed in
the clearest, deepest ice near the bottom centre of the main array,
lowering IceCube's effective energy threshold from $\sim$100~GeV down to a
few GeV -- the sub-array this document's entire analysis is built on
\parencite{Abbasi2023IceCube}. This document concerns the real 9-year
``golden event'' sample: well-reconstructed, predominantly track-like
(muon-neutrino charged-current) events, selected for their high purity and
good angular/energy reconstruction, used by the collaboration's own
atmospheric-oscillation analysis.

\begin{notebox}
\textbf{Scope of this document.} \texttt{tpeanuts} reproduces IceCube's own
real event-by-event weighted-Monte-Carlo reweighting technique, a real
3-flavour Earth-matter oscillation probability, the real Honda atmospheric
flux calculation, and the real per-bin detector-systematics hypersurface
correction, evaluated at the collaboration's own published best-fit
systematics point. It does \emph{not} re-fit the 5 detector-systematics
parameters as free nuisances (they are held fixed, a deliberate,
documented scope decision -- \cref{ch:detector-implementation}), nor
reproduce a full atmosphere-production-height treatment (production is
treated at the top of the atmosphere, a sub-percent-level simplification).
\Cref{ch:results} reports, quantitatively and without adjustment, what
these simplifications do and do not explain about a real discrepancy
between this document's own fitted muon-background normalization and
IceCube's own published value.
\end{notebox}

\section{Atmospheric neutrino production and the oscillation signal}

\begin{theorybox}
Atmospheric neutrinos are produced by cosmic-ray-induced particle showers
in the upper atmosphere: charged pions and kaons decay into muons and muon
neutrinos, and the muons themselves subsequently decay, producing further
muon and electron neutrinos -- a well-understood, precisely calculated
flux \parencite{Honda2015Flux} spanning many decades in energy and the
full range of zenith angles. Neutrinos arriving from below the horizon
(upgoing) have crossed some or all of the Earth's diameter since
production, at baselines ranging from a few hundred kilometres (near
horizontal) to $\sim$12{,}700~km (straight through the core) -- a natural,
continuously-varying baseline that, combined with the atmospheric flux's
own wide energy range, makes the upgoing $\nu_\mu$ disappearance pattern a
direct, high-statistics probe of $\theta_{23}$ and $\Delta m^2_{3l}$ (the
same oscillation minimum first resolved by Super-Kamiokande's own 1998
atmospheric-neutrino discovery, refined here with IceCube DeepCore's own
real GeV-scale event sample and a real Earth-matter treatment).
\end{theorybox}

\begin{notebox}
Downgoing (production directly overhead, negligible propagation baseline)
and near-horizontal events carry little oscillation information; this
document's real analysis binning (\cref{ch:data-release}) accordingly
spans $\cos\theta_z\in[-1,\,0.1]$ (upgoing to just past horizontal) only,
matching the real public release's own binning exactly.
\end{notebox}

\section{Event-by-event weighted Monte Carlo: why IceCube is different}
\label{sec:weighted-mc-intro}

\begin{theorybox}
Every other detector in this project (Daya Bay, Borexino, SNO) provides
either a continuous differential flux/cross section or an officially
extracted pseudo-observable, folded through a detector response using
\texttt{tpeanuts.detector.common.event\_rate}. IceCube's own public data
release provides neither: instead, it ships $\sim$400{,}000 individually
simulated Monte Carlo events, each carrying its own \emph{true} (simulated)
energy and direction, its \emph{reconstructed} analysis-bin location, and a
real physical weight (units GeV~cm$^2$~sr) that must be multiplied by the
atmospheric flux and a livetime to become a physical event count -- the
standard ``weighted Monte Carlo'' technique used throughout atmospheric-neutrino
oscillation analyses, and the reason \texttt{tpeanuts.detector.icecube}
reweights real simulated events directly (\cref{ch:detector-implementation})
rather than reusing the continuous-fold machinery every other detector
package in this project shares.
\end{theorybox}

\section{Data-taking period and physics goals}

\begin{theorybox}
The real public data release used throughout this document spans 9 years
of DeepCore operation, and reports IceCube's own real published best fit
(Table III/IV of \cite{Abbasi2023IceCube}, normal ordering, for later
comparison in \cref{ch:results}):
\begin{equation}
  \sin^2\theta_{23} = 0.51, \qquad
  \Delta m^2_{32} = 2.41^{+0.07}_{-0.07}\times10^{-3}\,{\rm eV}^2, \qquad
  \text{atmospheric-muon norm.\ scale} = 1.39,
\end{equation}
with $\theta_{23}$ measured via the atmospheric $\nu_\mu$ disappearance
minimum's own energy/zenith-angle position, and the muon-background
normalization scale a real, published nuisance-parameter result of the
collaboration's own joint fit (not an oscillation parameter, but a
detector/background quantity this document's own \cref{ch:results} audits
directly against its own fitted value).
\end{theorybox}
